\documentclass[10pt, a4paper]{article}

% --- Packages ---
\usepackage[top=2cm, bottom=2cm, left=2cm, right=2cm]{geometry}
\usepackage{graphicx} % For images
\usepackage{amsmath}  % For math/equations
\usepackage{hyperref} % For clickable links
\usepackage{caption}  % For styling captions
\usepackage{titlesec} % For controlling section title sizes
\usepackage{float}
\usepackage{xcolor} % Required for defining and using colors
\usepackage{hyperref} % Keep this as your last \usepackage
\usepackage{color}
\usepackage{soul}

% ברירת המחדל של המרקור היא צהוב. אם תרצי לשנות למשל לטורקיז/תכלת, הוסיפי גם את השורה הזו:
% \sethlcolor{cyan}

\hypersetup{
    colorlinks=true, % Removes the default boxes and colors the text instead
    linkcolor=blue,  % Color for internal links (e.g., Figure and Table references)
    citecolor=blue,  % Color for bibliographical citations
    urlcolor=blue    % Color for external web URLs
}
% --- Document Formatting ---
\setlength{\parindent}{0pt}
\setlength{\parskip}{0.5em}
\captionsetup{font=scriptsize, labelfont=bf}

% --- Customize Title Sizes ---
\titleformat{\section}
  {\large\bfseries}{\thesection}{1em}{}
\titleformat{\subsection}
  {\normalsize\bfseries}{\thesubsection}{1em}{}

\usepackage{fancyhdr}
\pagestyle{fancy}
\fancyhf{} % מנקה את הגדרות ברירת המחדל
\fancyhead[L]{\footnotesize Weekly Summary \#6 \today} % יופיע למעלה בשמאל
\fancyhead[R]{\footnotesize Rachel Broner $\vert$ Miri Adler's Lab} % יופיע למעלה בימין
\fancyfoot[C]{\thepage} % מספר עמוד למטה באמצע

% \usepackage{mathptmx}

\usepackage{setspace}
\onehalfspacing % For 1.5 line spacing


\begin{document}

\begin{center}
    {\Large \textbf{Weekly Summary \#6 $\vert$ Association Rule Mining}} \\[0.5em]
    {\large Part 1 $\vert$ Directional positive pairwise rules} \\[0.5em]
    \textit{GVHD Analysis} $\vert$ Week of: \today
\end{center}
\vspace{0.2cm}
\hrule
\vspace{0.5cm}

% ---------------------------------------------------------
\section{First biological analysis: directional positive pairwise rules}

\subsection{Q1 -- Which pairwise directional rules are shared across many FOVs/patients?}

\textbf{How this was calculated.}
The rules shown are those that rank among the 5 strongest positive pairwise rules
(two cell types, lift $>1$, no rule with the same cell type on both sides) of at least
one FOV, keeping the 30 found in the most FOVs/patients.
Each heatmap cell counts the FOVs / patients in which the rule was found; the percentage
is that count out of all FOVs / patients, or out of that stage's FOVs / patients in a
stage column.
Each bar is the rule's cell type as a fraction of all cells in an FOV, averaged over the
FOVs in which that rule was found (antecedent to the left, consequent to the right).


\begin{figure}[H]
    \centering
    \includegraphics[width=\textwidth]{summary_downloads/q1_prevalence_by_fov.pdf}
    \includegraphics[width=\textwidth]{summary_downloads/q1_prevalence_by_patient.pdf}
    \caption{Prevalence of the top pairwise directional rules across disease stages,
    counted by FOV (top) and by patient (bottom). Counting by biopsy gives an almost
    identical picture to counting by patient.}
    \label{fig:q1_prevalence}
\end{figure}

\subsection{Q2 -- Which rules are patient-specific or condition-specific?}

\textbf{How this was calculated}
Rules are selected as in Q1, run separately on the colon FOVs and on the duodenum FOVs.
Each heatmap cell counts the FOVs in which the rule was found; the percentage is that
count out of all the organ's FOVs, or out of that stage's FOVs in a stage column.
In the fingerprint, each cell is the share of that patient's FOVs in which the rule was
found; patients and rules are ordered so that similar ones sit next to each other, and the
strips above give each patient's stage.
Stage is the \emph{pathological} score throughout; Control and Control\_S are the FOVs with
no score, i.e.\ the controls.

\begin{figure}[H]
    \centering
    \includegraphics[width=\textwidth]{summary_downloads/q2_prevalence_colon.pdf}
    \includegraphics[width=\textwidth]{summary_downloads/q2_prevalence_duodenum.pdf}
    \caption{Top rules in the colon (top) and in the duodenum (bottom). Every Control\_S
    FOV is duodenum, so the colon panel has no Control\_S column.}
    \label{fig:q2_organ}
\end{figure}

\begin{figure}[H]
    \centering
    \includegraphics[width=\textwidth]{summary_downloads/q2_patient_fingerprint.pdf}
    \caption{Which rules each patient has, for the most widespread rules. Each column is
    one patient.}
    \label{fig:q2_fingerprint}
\end{figure}

\begin{figure}[H]
    \centering
    \includegraphics[width=\textwidth]{summary_downloads/q2_patient_fingerprint_colon.pdf}
    \caption{The same picture inside the colon only. The rule set is chosen within the organ,
    so the rows differ from the figure above.}
    \label{fig:q2_fingerprint_colon}
\end{figure}

\begin{figure}[H]
    \centering
    \includegraphics[width=\textwidth]{summary_downloads/q2_patient_fingerprint_duodenum.pdf}
    \caption{The same picture inside the duodenum only.}
    \label{fig:q2_fingerprint_duodenum}
\end{figure}

% ---------------------------------------------------------
\subsection{Q3 -- Which rules have high confidence/conviction/lift and also appear in more
than one FOV?}

\textbf{How this was calculated.}
No top-$N$ selection is used here: every positive pairwise rule is kept, one row each.
A rule's lift, confidence and conviction are the median of its values over the FOVs in
which it was found.
Of these, only the rules found in at least two patients are shown -- a stricter test
than the two FOVs asked for.
Each dot is one rule: its position gives lift and conviction, both on a log scale, its size
the number of patients, and its colour the confidence.
Two rules reach infinite conviction (confidence $=1$) and are drawn at the top of the plot
with a red ring.

\begin{figure}[H]
    \centering
    \includegraphics[width=0.85\textwidth]{summary_downloads/q3_lift_vs_conviction.pdf}
    \caption{Every positive pairwise rule found in at least two patients, by median lift
    and median conviction. Only the strongest few are named.}
    \label{fig:q3_strength}
\end{figure}

% ---------------------------------------------------------
\subsection{Q4 -- Which rules may simply reflect cell-type abundance?}

\textbf{How this was calculated.}
Every positive pairwise rule is used, with no top-$N$ selection and nothing averaged:
each dot is one rule in one FOV.
The horizontal axis is the share of cells in that FOV that are the rule's antecedent (left
panel) or consequent (right panel); the vertical axis is that rule's lift, confidence or
conviction in the same FOV.
The thick line is the median per bin, a guide to the direction only.
Infinite conviction is drawn at the top of the panel with a red ring rather than dropped.
The last figure is the Q1 patient heatmap re-ordered so that the rules with the most common
cell types sit at the top.

\begin{figure}[H]
    \centering
    \includegraphics[width=0.9\textwidth]{summary_downloads/q4_abundance_lift.pdf}
    \includegraphics[width=0.9\textwidth]{summary_downloads/q4_abundance_confidence.pdf}
    \includegraphics[width=0.9\textwidth]{summary_downloads/q4_abundance_conviction.pdf}
    \caption{Each score against how common the rule's cell types are, one dot per rule
    per FOV.}
    \label{fig:q4_abundance}
\end{figure}

\begin{figure}[H]
    \centering
    \includegraphics[width=\textwidth]{summary_downloads/q4_prevalence_by_abundance.pdf}
    \caption{The top rules by patient, re-ordered so that the rules with the most common
    cell types sit at the top.}
    \label{fig:q4_sorted}
\end{figure}

% ---------------------------------------------------------

\end{document}
